\begin{frame}[fragile]{PKI}

\begin{itemize}
	\item Just as social infrastructure maintains the security and authenticity of money, public-key infrastructure (PKI) guarantees information security and authenticity for public-key encryption.
	\item In other words, because of PKI, we can use public-key encryption to exchange emails, do business over the internet, and perform other actions with peace of mind.
	\item The KDC (Key Distribution Center) is an example of PKI.
\end{itemize}

\end{frame}

\begin{frame}[fragile]{}

\begin{itemize}
	\item Anyone can register public keys in KDC, or  KDC can create public keys whenever necessary. 
	\item In this example, user A register his public key and get a certificate. 
	\item When user B needs a user A's public key, she can access KDC to get A's authenticated public keys. 
\end{itemize}

\begin{center}
\includegraphics[width=0.5\textwidth]{topics/SE_tools2/security/pic/kdc/kdc.png}
\end{center} 
\end{frame}

\begin{frame}[fragile]{}

\begin{itemize}
	\item Using his own private key, User A attaches a digital signature he has generated to the message. Then he sends the message to User B.
	\item Using user A's public key, User B verifies the digital signature on the received message. If the decryption of the digital signature with User A's public key matches the message sent, the message is legitimate.
\end{itemize}

\begin{columns}%
\begin{column}{0.5\textwidth}%
\begin{center}%
\includegraphics[width=0.99\textwidth]{topics/SE_tools2/security/pic/kdc/sender.png}
\end{center}%
\end{column}%
$\rightarrow$
\begin{column}{0.5\textwidth}%
\begin{center}%
\includegraphics[width=0.99\textwidth]{topics/SE_tools2/security/pic/kdc/receiver.png}
\end{center}%
\end{column}%
\end{columns}%

\end{frame}

\begin{frame}[fragile]{Kerberos KDC Protocol}

\begin{itemize}
	\item Kerberos is the most widely used KDC protocol. 
	\item When we need to provide KDC service, we need to install Kerberos server. 
\end{itemize}

\begin{columns}%
\begin{column}{0.5\textwidth}%
\begin{center}%
\includegraphics[width=0.99\textwidth]{topics/SE_tools2/security/pic/kdc/k1.png}
\end{center}%
\end{column}%
\begin{column}{0.5\textwidth}%
\begin{center}%

\includegraphics[width=0.5\textwidth]{topics/SE_tools2/security/pic/kdc/k2.png}
\end{center}%
\end{column}%
\end{columns}%

\end{frame}

\begin{frame}[fragile]{The Trusted Certificate Authority (CA)}

\begin{itemize}
	\item How do we know this public key belongs to user A?
	\item So, we need the trusted CA. 
	\item CA certifies and publishes user A's request.
	\item In this case, user A requests A's public key authorization. 
\end{itemize}
\end{frame}

\begin{frame}[fragile]{CA Certifies A's Request}

\begin{itemize}
	\item When the decryption of the received application and the public key are verified, the certificate authority will grant user A the certificate.
\end{itemize}

\begin{center}
\includegraphics[width=0.8\textwidth]{topics/SE_tools2/security/pic/kdc/k7.png}
\end{center} 
\end{frame}

\begin{frame}[fragile]{CA, Users A and B to share the Public key}

\begin{itemize}
	\item User A registers their public key with the CA, which publishes a certificate.
	\item The certificate is encrypted with the CA's private key, so we need CA's public key to encrypt the certificate. 
\end{itemize}

\begin{columns}%
\begin{column}{0.33\textwidth}%
\begin{center}%
\includegraphics[width=0.99\textwidth]{topics/SE_tools2/security/pic/kdc/k3.png}
\end{center}%
\end{column}%
\begin{column}{0.33\textwidth}%
\begin{center}%
\includegraphics[width=0.99\textwidth]{topics/SE_tools2/security/pic/kdc/k4.png}
\end{center}%
\end{column}%
\end{columns} 
\end{frame}

\begin{frame}[fragile]{}

\begin{itemize}
	\item User B downloads User A's certificate from the repository.
\end{itemize}

\begin{columns}%
\begin{column}{0.43\textwidth}%
\begin{center}%
\includegraphics[width=0.99\textwidth]{topics/SE_tools2/security/pic/kdc/k4.png}
\end{center}%
\end{column}%
\begin{column}{0.43\textwidth}%
\begin{center}%
\includegraphics[width=1.1\textwidth]{topics/SE_tools2/security/pic/kdc/k5.png}
\end{center}% 
\end{column}%
\end{columns}%
\end{frame}

\begin{frame}[fragile]{}

\begin{itemize}
	\item User B decrypts user A's certificate using the CA's public key.
	\item User B now has a public key in the certificate and another public key that is certified by CA. 
	\item User B then compares the public key and digital signature key. 
	\item If the two keys are the same, User B has verified that the message came from User A.
\end{itemize}

\begin{center}
\includegraphics[width=0.5\textwidth]{topics/SE_tools2/security/pic/kdc/k6.png}
\end{center} 
\end{frame}

\begin{frame}[fragile]{Automatization of the CA related Process}

\begin{columns}%
\begin{column}{0.7\textwidth}%

\begin{itemize}
  \item In practice, all of the validation and authentication process is automatically processed by software tools. 
  \item Web browsers, private software registration cards, and card readers, all have the information about the CAs to automatize the whole process. 
\end{itemize}

\end{column}%
\begin{column}{0.3\textwidth}%
\begin{center}%
\includegraphics[width=0.9\textwidth]{topics/SE_tools2/security/pic/kdc/cert.png}  
\end{center}%
\end{column}%
\end{columns}%
\end{frame}
